% Generado automáticamente desde notas/09_graficos_propios.md con md2tex.py; no editar a mano.
\begin{nota}
Estado: COMPLETA (29/09/2026).
\end{nota}

\section{Resumen}

\begin{itemize}
\item La web tiene \textbf{17 gráficos o mapas} en 16 entradas de inventario (13 ECharts, 1 violín SVG propio, 1 mapa mundial d3, 1 trazado SVG y 1 matriz-tabla), \textbf{15 tarjetas KPI} y unas \textbf{25 tablas}. Todos los gráficos van dentro de \code{Chart\allowbreak{}Figure}, con título, resumen en texto y la tabla de datos desplegable. La base de accesibilidad es buena y mejor que la de la mayoría de webs de F1.
\item Réplica del TFG: las páginas de las figuras 5.23-5.33 están todas (mapa, récords de pilotos y constructores, temporada, clasificación, vuelta a vuelta, neumáticos, tiempos, ritmo, boxes y detalle de piloto). Además hay extras: evolución del campeonato, telemetría con trazado, fichas de constructor y calidad de datos.
\item Problemas principales:
\begin{enumerate}
\item \textbf{Color sin identidad}: la paleta Okabe-Ito de 8 colores se repite en carreras de 20 pilotos. Los compañeros no comparten color de equipo y en el violín dos pilotos distintos salen con el mismo color sin otra diferencia.
\item \textbf{Solapamientos}: stints y matriz de neumáticos cuentan lo mismo. Hay tres gráficos y dos tablas de compañeros con la misma información. Las tarjetas de temporada repiten la tabla.
\item \textbf{Errores de lectura concretos}: la «referencia» de clasificación no es la pole. La duración de parada es el tiempo en el pit lane, no la parada. Las temporadas sin carrera desaparecen del eje. Ferrari tiene barras a 0 en 1950-57. Las etiquetas del final de la evolución del campeonato se solapan. «1,000 m» aparece con el formato inglés.
\item \textbf{Huecos de relato}: no hay página de circuito, ni «gap al líder», ni comparador libre de pilotos, ni fiabilidad o edades. Tampoco se ve cuándo se decidió el título.
\end{enumerate}
\item Propongo 26 visualizaciones nuevas (P1-P26, más la P1b), todas comprobadas contra \code{data/\allowbreak{}gold/\allowbreak{}f1.\allowbreak{}duckdb}, y una navegación con dos entradas nuevas: \textbf{Circuitos} y \textbf{Comparar}. Además, \textbf{Historia} agruparía Récords y las vistas por eras.
\end{itemize}

Fuentes: código de \code{web/\allowbreak{}src/\allowbreak{}app/\allowbreak{}[lang]/\allowbreak{}**} y \code{web/\allowbreak{}src/\allowbreak{}components/\allowbreak{}**} (sin modificar), la web desplegada (navegada el 29/09/2026 a 640 px y a 375 px de ancho), el TFG (texto de las figuras 5.23-5.33 extraído con \code{pdftotext}, porque el PDF no deja leerlo directamente) y \code{data/\allowbreak{}gold/\allowbreak{}f1.\allowbreak{}duckdb} en modo de solo lectura.

Capturas en \code{docs/\allowbreak{}informe\_\allowbreak{}situacion/\allowbreak{}figuras/\allowbreak{}}:

{\small\sloppy\setlength{\tabcolsep}{5pt}\renewcommand{\arraystretch}{1.15}
\begin{longtable}{@{}>{\raggedright\arraybackslash}p{0.393\dimexpr(\linewidth-20pt)\relax}>{\raggedright\arraybackslash}p{0.607\dimexpr(\linewidth-20pt)\relax}@{}}
\toprule
\textbf{Fichero} & \textbf{Contenido} \\
\midrule\endhead
\bottomrule\endlastfoot
\code{web\_\allowbreak{}lap\_\allowbreak{}chart\_\allowbreak{}bahrain2024.\allowbreak{}png} & Posiciones vuelta a vuelta, Bahrain 2024 \\
\code{web\_\allowbreak{}tiempos\_\allowbreak{}vuelta\_\allowbreak{}bahrain2024.\allowbreak{}png} & Tiempos por vuelta (20 pilotos) \\
\code{web\_\allowbreak{}violin\_\allowbreak{}bahrain2024.\allowbreak{}png} & Violines de ritmo \\
\code{web\_\allowbreak{}stints\_\allowbreak{}bahrain2024.\allowbreak{}png} & Estrategia de neumáticos \\
\code{web\_\allowbreak{}quali\_\allowbreak{}gap\_\allowbreak{}bahrain2024.\allowbreak{}png} & \% sobre el mejor tiempo (etiquetas solapadas) \\
\code{web\_\allowbreak{}telemetria\_\allowbreak{}bahrain2024.\allowbreak{}png} & Velocidad, acelerador y freno \\
\code{web\_\allowbreak{}temporada2021\_\allowbreak{}kpi\_\allowbreak{}progresion.\allowbreak{}png} & Tarjetas KPI y evolución del campeonato de 2021 (etiquetas solapadas) \\
\code{web\_\allowbreak{}piloto\_\allowbreak{}schumacher\_\allowbreak{}kpi\_\allowbreak{}puntos.\allowbreak{}png} & KPIs y puntos por temporada (los años 2007-2009 desaparecen del eje) \\
\code{web\_\allowbreak{}piloto\_\allowbreak{}schumacher\_\allowbreak{}companeros.\allowbreak{}png} & Gráficos de compañeros \\
\end{longtable}}

\medskip\noindent\textcolor{muted!40}{\rule{\linewidth}{0.3pt}}\medskip

\section{Método}

\begin{enumerate}
\item Lectura completa de las 17 páginas y de los 17 componentes visuales: tipos de marca, escalas, colores, datos de entrada y textos alternativos.
\item Navegación por la web desplegada: \code{/\allowbreak{}es}, \code{/\allowbreak{}es/\allowbreak{}seasons/\allowbreak{}2021}, \code{/\allowbreak{}es/\allowbreak{}races/\allowbreak{}1102} (Bahrain 2024) y sus pestañas, \code{/\allowbreak{}es/\allowbreak{}drivers/\allowbreak{}michael-\allowbreak{}schumacher}, \code{/\allowbreak{}es/\allowbreak{}constructors/\allowbreak{}ferrari} y \code{/\allowbreak{}es/\allowbreak{}records}. Se probó a 640 px y a 375 px de ancho.
\item Consultas de cobertura en DuckDB para saber qué gráficos se pueden dibujar y desde qué año (sección 5.0).
\end{enumerate}

\medskip\noindent\textcolor{muted!40}{\rule{\linewidth}{0.3pt}}\medskip

\section{Inventario completo}

Leyenda de accesibilidad:

\begin{itemize}
\item \textbf{F} = \code{Chart\allowbreak{}Figure}: título, resumen textual y tabla desplegable.
\item \textbf{img} = contenedor \code{role=\allowbreak{}"img"} con \code{aria-\allowbreak{}label}.
\item \textbf{T} = tabla semántica: \code{caption}, \code{scope} y región desplazable con el teclado.
\end{itemize}

\subsection{Gráficos y mapas}

{\scriptsize\sloppy\setlength{\tabcolsep}{4pt}\renewcommand{\arraystretch}{1.15}
\begin{longtable}{@{}>{\raggedright\arraybackslash}p{0.015\dimexpr(\linewidth-72pt)\relax}>{\raggedright\arraybackslash}p{0.094\dimexpr(\linewidth-72pt)\relax}>{\raggedright\arraybackslash}p{0.112\dimexpr(\linewidth-72pt)\relax}>{\raggedright\arraybackslash}p{0.152\dimexpr(\linewidth-72pt)\relax}>{\raggedright\arraybackslash}p{0.152\dimexpr(\linewidth-72pt)\relax}>{\raggedright\arraybackslash}p{0.139\dimexpr(\linewidth-72pt)\relax}>{\raggedright\arraybackslash}p{0.152\dimexpr(\linewidth-72pt)\relax}>{\raggedright\arraybackslash}p{0.139\dimexpr(\linewidth-72pt)\relax}>{\raggedright\arraybackslash}p{0.044\dimexpr(\linewidth-72pt)\relax}@{}}
\toprule
\textbf{\#} & \textbf{Ruta} & \textbf{Componente} & \textbf{Tipo} & \textbf{Datos (endpoint)} & \textbf{Pregunta que responde} & \textbf{Interacción} & \textbf{Accesibilidad} & \textbf{TFG} \\
\midrule\endhead
\bottomrule\endlastfoot
G1 & \code{/\allowbreak{}} (\#mapa) & \code{World\allowbreak{}Map} (d3-geo, SVG en servidor, sin JS) & Mapa de símbolos proporcionales (área ∝ √GP), proyección Equal Earth & \code{/\allowbreak{}circuits?season\_\allowbreak{}from\&season\_\allowbreak{}to}; en la vista por país, centroide ponderado por carreras & ¿Dónde se ha corrido la F1 y cuánto? & Por país o por circuito (enlaces); rango de temporadas (formulario GET); \code{<title>} en cada círculo & \code{<title>}/\code{<desc>}, figcaption y tabla por país o circuito & 5.23 \\
G2 & \code{/\allowbreak{}seasons/\allowbreak{}[year]} & \code{Progression\allowbreak{}Chart} & Líneas de puntos acumulados por ronda (top 10), con etiqueta al final & \code{/\allowbreak{}seasons/\allowbreak{}\{y\}/\allowbreak{}standings/\allowbreak{}progression?top=\allowbreak{}10} & ¿Cómo se fraguó el campeonato? & Tooltip por eje; resaltado de la serie al pasar & F, img, tabla piloto × ronda & Nuevo \\
G3 & \code{/\allowbreak{}races/\allowbreak{}[id]/\allowbreak{}qualifying} & \code{Bar\allowbreak{}Chart} vertical & Barras del \% sobre el mejor tiempo de la sesión; la mejor en morado & \code{/\allowbreak{}races/\allowbreak{}\{id\}/\allowbreak{}qualifying?session=\allowbreak{}} → \code{gap\_\allowbreak{}to\_\allowbreak{}pole\_\allowbreak{}pct} & ¿A cuánto quedó cada piloto o equipo del mejor tiempo? & Pilotos o constructores; Q o Q-sprint & F, img & 5.27 \\
G4 & \code{/\allowbreak{}races/\allowbreak{}[id]/\allowbreak{}lap-\allowbreak{}chart} & \code{Lap\allowbreak{}Chart} & Bump chart de posición por vuelta, con parrilla en la vuelta 0 y código del piloto al final & \code{/\allowbreak{}races/\allowbreak{}\{id\}/\allowbreak{}laps} + \code{/\allowbreak{}results} (parrilla) & ¿Cómo evolucionó el orden en carrera? & Resaltado al pasar; tooltip por punto & F, img, tabla salida/llegada/mejor/vueltas lideradas & 5.28 \\
G5 & \code{/\allowbreak{}races/\allowbreak{}[id]/\allowbreak{}tyres} & \code{Stint\allowbreak{}Chart} & Barras horizontales apiladas por stint; color Pirelli + letra & \code{/\allowbreak{}races/\allowbreak{}\{id\}/\allowbreak{}stints} & ¿Qué estrategia hizo cada uno? & Tooltip & F, img, leyenda con letra, tabla & 5.29 \\
G6 & \code{/\allowbreak{}races/\allowbreak{}[id]/\allowbreak{}tyres} (\#matriz) & \code{Compound\allowbreak{}Matrix} & Tabla-mapa de calor piloto × vuelta (compuesto; vuelta de entrada recuadrada) & \code{/\allowbreak{}races/\allowbreak{}\{id\}/\allowbreak{}laps} & ¿Qué neumático llevaba en cada vuelta y cuándo paró? & \code{title} por celda; desplazamiento horizontal & T con texto \code{sr-\allowbreak{}only} por celda & 5.29 (matriz) \\
G7 & \code{/\allowbreak{}races/\allowbreak{}[id]/\allowbreak{}pace} (\#tiempos) & \code{Lap\allowbreak{}Times\allowbreak{}Chart} & Líneas de tiempo por vuelta; 2 dataZoom (vueltas y tiempos) & \code{/\allowbreak{}races/\allowbreak{}\{id\}/\allowbreak{}laps} & ¿Cómo evolucionó el ritmo de cada piloto? & Casillas por piloto, Todos/Ninguno, ocultar vueltas neutralizadas, zoom & F, img, tabla vuelta × piloto & 5.30 \\
G8 & \code{/\allowbreak{}races/\allowbreak{}[id]/\allowbreak{}pace} (\#ritmo) & \code{Violin\allowbreak{}Chart} (SVG propio, KDE de Silverman) & Violín + caja P25-P75 + mediana por piloto, ordenado por mediana & Laps filtrados: sin vuelta 1, boxes, SC/VSC/roja ni vueltas >150 \% de la mediana; mín. 5 vueltas & ¿Quién tuvo mejor ritmo de verdad y lo mantuvo? & Casillas por piloto; \code{<title>} por violín & F, img, tabla mediana/mejor/P25-P75/nº & 5.31 \\
G9 & \code{/\allowbreak{}races/\allowbreak{}[id]/\allowbreak{}pitstops} & \code{Bar\allowbreak{}Chart} vertical & Barras del tiempo medio por equipo (paradas ≤ 60 s); la mejor en morado & \code{/\allowbreak{}races/\allowbreak{}\{id\}/\allowbreak{}pitstops} & ¿Qué equipo perdió menos tiempo en boxes? & Tooltip & F, img & 5.32 \\
G10 & \code{/\allowbreak{}races/\allowbreak{}[id]/\allowbreak{}telemetry} (desde 2024) & \code{Telemetry\allowbreak{}Chart} & 3 paneles con el eje X compartido: velocidad, acelerador y freno frente a la distancia & \code{/\allowbreak{}races/\allowbreak{}\{id\}/\allowbreak{}telemetry?drivers=\allowbreak{}a,\allowbreak{}b} (vuelta más rápida de clasificación) & ¿Dónde gana tiempo un piloto frente a otro? & Selector A/B (formulario GET); eje enlazado & F, img, tabla tiempo/v. máx./v. mín./\% a fondo & Nuevo \\
G11 & \code{/\allowbreak{}races/\allowbreak{}[id]/\allowbreak{}telemetry} & \code{Track\allowbreak{}Map} & Trazado x/y coloreado por velocidad (azul→naranja), solo del piloto A & Igual que G10 & ¿Cómo es el circuito y dónde es rápido o lento? & Ninguna & \code{<title>}/\code{<desc>}; leyenda con \code{aria-\allowbreak{}hidden} & Nuevo \\
G12 & \code{/\allowbreak{}drivers/\allowbreak{}[id]} & \code{Season\allowbreak{}Points\allowbreak{}Chart} & Barras de puntos por temporada; título en morado y posición final encima & \code{/\allowbreak{}drivers/\allowbreak{}\{id\}/\allowbreak{}seasons} & ¿Cómo fue su carrera deportiva año a año? & Tooltip & F, img, tabla temporada/equipo/pos./pts/V/P/poles & 5.33 (parcial) \\
G13 & \code{/\allowbreak{}drivers/\allowbreak{}[id]} & \code{Teammate\allowbreak{}Points\allowbreak{}Chart} & Barras agrupadas: puntos propios frente a los de sus compañeros (trama diagonal) & \code{/\allowbreak{}drivers/\allowbreak{}\{id\}/\allowbreak{}teammates} agregado por temporada & ¿Sumó más que sus compañeros? & Leyenda y tooltip & F, img, textura además del color & 5.33 \\
G14 & \code{/\allowbreak{}drivers/\allowbreak{}[id]} & \code{Teammate\allowbreak{}Share\allowbreak{}Chart} ×2 & Barras apiladas al 100 \%: \% por delante en clasificación y en carrera & Igual & ¿Batió a sus compañeros en clasificación y en carrera? & Leyenda y tooltip & F, img, textura & 5.33 \\
G15 & \code{/\allowbreak{}constructors/\allowbreak{}[id]} & \code{Season\allowbreak{}Points\allowbreak{}Chart} & Igual que G12 & \code{/\allowbreak{}constructors/\allowbreak{}\{id\}/\allowbreak{}seasons} & Evolución del equipo & Tooltip & F, img & Nuevo \\
G16 & \code{/\allowbreak{}records} & \code{Bar\allowbreak{}Chart} horizontal & Títulos por piloto o constructor en el rango de años & \code{/\allowbreak{}rankings/\allowbreak{}\{entity\}?order\_\allowbreak{}by=\allowbreak{}championships} & ¿Quién tiene más títulos en esa época? & Entidad y rango de años & F, img & 5.24-5.25 \\
\end{longtable}}

\subsection{Tarjetas KPI}

{\footnotesize\sloppy\setlength{\tabcolsep}{5pt}\renewcommand{\arraystretch}{1.15}
\begin{longtable}{@{}>{\raggedright\arraybackslash}p{0.156\dimexpr(\linewidth-40pt)\relax}>{\raggedright\arraybackslash}p{0.134\dimexpr(\linewidth-40pt)\relax}>{\raggedright\arraybackslash}p{0.435\dimexpr(\linewidth-40pt)\relax}>{\raggedright\arraybackslash}p{0.275\dimexpr(\linewidth-40pt)\relax}@{}}
\toprule
\textbf{Ruta} & \textbf{Componente} & \textbf{Contenido} & \textbf{TFG} \\
\midrule\endhead
\bottomrule\endlastfoot
\code{/\allowbreak{}seasons/\allowbreak{}[year]} & \code{Highlight} ×3 & Campeón (o líder), equipo campeón (o líder) y victorias del campeón & 5.26 (tarjetas de campeón y equipo) \\
\code{/\allowbreak{}drivers/\allowbreak{}[id]} & \code{Stat} ×6 & Títulos, victorias, podios, poles, vueltas rápidas y salidas & 5.33 (tarjetas superiores) \\
\code{/\allowbreak{}constructors/\allowbreak{}[id]} & \code{Stat} ×6 & Títulos, victorias, podios, poles, dobletes y carreras & Nuevo \\
\code{/\allowbreak{}quality} & Párrafo resumen & «N controles superados, M fallidos» & Nuevo \\
\end{longtable}}

\subsection{Tablas}

{\small\sloppy\setlength{\tabcolsep}{5pt}\renewcommand{\arraystretch}{1.15}
\begin{longtable}{@{}>{\raggedright\arraybackslash}p{0.181\dimexpr(\linewidth-30pt)\relax}>{\raggedright\arraybackslash}p{0.409\dimexpr(\linewidth-30pt)\relax}>{\raggedright\arraybackslash}p{0.409\dimexpr(\linewidth-30pt)\relax}@{}}
\toprule
\textbf{Ruta} & \textbf{Tabla} & \textbf{Observación} \\
\midrule\endhead
\bottomrule\endlastfoot
\code{/\allowbreak{}} & Top 10 de la última carrera; top 5 de pilotos y de constructores & Puerta de entrada; repite la página de temporada \\
\code{/\allowbreak{}seasons} & Tarjetas por década con los campeones & Es una lista de enlaces, sin datos \\
\code{/\allowbreak{}seasons/\allowbreak{}[year]} & Pilotos, constructores y calendario con ganador & Réplica de la figura 5.26 \\
\code{/\allowbreak{}races/\allowbreak{}[id]} & Resultado de carrera o sprint & Pills de vuelta rápida y dato corregido \\
\code{/\allowbreak{}races/\allowbreak{}[id]/\allowbreak{}qualifying} & Q1/Q2/Q3 con el mejor segmento en morado; \% & La columna de \% repite G3 \\
\code{/\allowbreak{}races/\allowbreak{}[id]/\allowbreak{}pitstops} & Lista de paradas; pasos por el pit lane que no son paradas & La segunda tabla es original y valiosa (SC, bandera roja, sanciones) \\
\code{/\allowbreak{}drivers}, \code{/\allowbreak{}constructors} & Top 50 por victorias o búsqueda & Sin orden por columnas (a diferencia de \code{/\allowbreak{}records}) \\
\code{/\allowbreak{}drivers/\allowbreak{}[id]} & Detalle de compañeros (temporada × compañero) & Más la tabla por temporada, repetida en los 3 gráficos \\
\code{/\allowbreak{}records} & Ranking ordenable (7 métricas) con rango de años & Réplica de las figuras 5.24-5.25 \\
\code{/\allowbreak{}quality} & Controles con umbral y controles informativos & Página original del proyecto \\
\end{longtable}}

\medskip\noindent\textcolor{muted!40}{\rule{\linewidth}{0.3pt}}\medskip

\section{Evaluación crítica}

\subsection{Problemas transversales}

\begin{enumerate}
\item \textbf{Color y paleta.}
\begin{itemize}
\item \code{LIGHT\_\allowbreak{}SERIES} tiene 8 colores. A partir del 9.º piloto el color se repite y cambia el trazo (discontinuo o punteado), lo que funciona en líneas finas.
\item En el \textbf{violín} (G8) no hay trazo que valga: Hülkenberg y Piastri salen del mismo marrón, y Stroll y Pérez del mismo naranja (captura \code{web\_\allowbreak{}violin\_\allowbreak{}bahrain2024.\allowbreak{}png}).
\item No se usa el color del equipo. El aficionado espera Ferrari en rojo y los dos compañeros con el mismo tono. \textbf{Dato que falta}: no hay colores de equipo en gold. \code{stg\_\allowbreak{}fastf1\_\allowbreak{}\_\allowbreak{}results.\allowbreak{}team\_\allowbreak{}color} existe, pero su vista apunta a un bronze que no está disponible (da error al consultarla), y además solo cubre 2018+. Propuesta: seed \code{constructor\_\allowbreak{}colors.\allowbreak{}csv} (id, temporada desde/hasta, hex claro y oscuro) y usar variantes claras y oscuras para distinguir a los compañeros, manteniendo el trazo o la forma.
\end{itemize}
\item \textbf{Ejes de temporadas como categorías.} \code{Season\allowbreak{}Points\allowbreak{}Chart} y \code{Teammate\allowbreak{}Points\allowbreak{}Chart} usan \code{type:\allowbreak{} "category"}, así que las temporadas sin carrera desaparecen. En Schumacher, 2006 va pegado a 2010 y parece que no hubo retirada. Debe ser un eje de valor (años) o marcar los huecos.
\item \textbf{Puntos entre eras.} Los gráficos de puntos por temporada (G12, G15) mezclan sistemas de puntuación: 9 puntos por victoria en 1950 frente a 25 hoy, y ahora más carreras. Los 148 puntos de Schumacher en 2004 parecen poco frente a los 575 de Verstappen en 2023. Hace falta una métrica normalizada (sección 5, P14).
\item \textbf{Base del eje en cero donde no aporta.}
\begin{itemize}
\item En la media de paradas (G9) las barras van de 24,1 a 27,2 s y el eje de 0 a 30: todas parecen iguales. En barras hay que empezar en 0, así que es mejor un gráfico de puntos o \emph{dot plot} con el eje ajustado, o mostrar la diferencia con el mejor.
\item En telemetría (G10) la velocidad empieza en 0 km/h y la mitad inferior del panel queda vacía.
\end{itemize}
\item \textbf{Formato numérico en ECharts.} Los ejes usan el formato por defecto de ECharts y aparece «1,000 m» en la página española. Los tooltips sí usan \code{to\allowbreak{}Locale\allowbreak{}String}. Falta un \code{axis\allowbreak{}Label.\allowbreak{}formatter} con \code{Intl.\allowbreak{}Number\allowbreak{}Format(\allowbreak{}locale)}.
\item \textbf{Etiquetas que se solapan.}
\begin{itemize}
\item \code{Progression\allowbreak{}Chart}: \code{label\allowbreak{}Layout.\allowbreak{}move\allowbreak{}Overlap} no se aplica a \code{end\allowbreak{}Label}. En 2021 «Max Verstappen» y «Lewis Hamilton» se pisan, y «Carlos Sainz Jr./Lando Norris/Charles Leclerc» no se leen.
\item Barras verticales de G3 con 20 pilotos: los valores «1,78 \%», «1,79 \%»… chocan entre sí y los nombres van girados 40°.
\end{itemize}
\item \textbf{Móvil (375 px).}
\begin{itemize}
\item El lap chart se comprime en \textasciitilde{}250 px de ancho y los intercambios en boxes forman una madeja ilegible.
\item Las pestañas de carrera y la navegación principal se desplazan en horizontal con una barra visible (se ve «Vuelta a vuelta» y la siguiente pestaña cortada).
\item El violín obliga a desplazarse en horizontal.
\item El mapa del inicio funciona bien porque es un SVG que se escala.
\end{itemize}
\item \textbf{Sin enlaces entre gráficos.} Ningún gráfico enlaza a la ficha del piloto o de la carrera (ECharts permite \code{on(\allowbreak{}'click')}) ni recuerda la selección de pilotos entre pestañas. Las casillas de G7 y G8 son independientes: el aficionado marca 3 pilotos en tiempos y tiene que volver a marcarlos en el violín.
\item \textbf{Accesibilidad.} Lo correcto: figura, resumen, tabla, letras en compuestos, texturas en compañeros y \code{prefers-\allowbreak{}reduced-\allowbreak{}motion}. Mejorable:
\begin{itemize}
\item \code{aria.\allowbreak{}enabled:\allowbreak{} false} es deliberado (la tabla sustituye al gráfico), pero el tooltip no es accesible con el teclado. No hay forma de «recorrer» la serie.
\item El texto del violín es de 10 px dentro de un \code{view\allowbreak{}Box} que se escala hacia abajo: en la práctica queda por debajo de 8 px (captura).
\item El resumen de G16 dice «Lewis Hamilton suma más títulos: 7» cuando empata con Schumacher. Los resúmenes deberían detectar empates.
\end{itemize}
\item \textbf{Rendimiento y carga.} Cada pestaña de carrera vuelve a pedir \code{/\allowbreak{}laps} (unas 1 200 filas) y \code{/\allowbreak{}results}. Hay \code{cache} por petición, pero no entre pestañas. Es aceptable, aunque con Render free la primera visita tarda.
\end{enumerate}

\subsection{Evaluación de cada elemento}

{\footnotesize\sloppy\setlength{\tabcolsep}{5pt}\renewcommand{\arraystretch}{1.15}
\begin{longtable}{@{}>{\raggedright\arraybackslash}p{0.155\dimexpr(\linewidth-50pt)\relax}>{\raggedright\arraybackslash}p{0.159\dimexpr(\linewidth-50pt)\relax}>{\raggedright\arraybackslash}p{0.229\dimexpr(\linewidth-50pt)\relax}>{\raggedright\arraybackslash}p{0.229\dimexpr(\linewidth-50pt)\relax}>{\raggedright\arraybackslash}p{0.229\dimexpr(\linewidth-50pt)\relax}@{}}
\toprule
\textbf{\#} & \textbf{Valoración} & \textbf{Puntos fuertes} & \textbf{Problemas concretos} & \textbf{Recomendación} \\
\midrule\endhead
\bottomrule\endlastfoot
G1 Mapa & \textbf{Útil} & Sin JS; rango de años; 2 vistas; tabla & Europa es una mancha de círculos que se solapan; el centroide por país sitúa a EE. UU. en medio del país y no en sus circuitos; no muestra el tiempo & Añadir un zoom a Europa (inserto) o una vista de «calendario por país × década» (P9); color por última temporada (vigente o histórico) \\
G2 Evolución campeonato & \textbf{Muy útil} & Cuenta la temporada; tabla completa & Etiquetas solapadas; 10 líneas con colores repetidos; los puntos absolutos esconden la diferencia en la lucha por el título & Añadir el modo «diferencia con el líder» (P1), resaltar por defecto los 2-3 aspirantes, marcar la ronda en que se decidió (\code{drivers\_\allowbreak{}championship\_\allowbreak{}decider}) \\
G3 \% sobre el mejor tiempo & \textbf{Útil, con un fallo de relato} & Métrica fiel al TFG; modo por constructores & (1) La referencia es el mejor tiempo de cualquier segmento: en Bahrain 2024 la «referencia» es Leclerc (Q2, 1:29.165) y Verstappen, que hizo la pole, aparece a +0,016 \%; el resumen dice «Leclerc marcó la referencia», lo que confunde. (2) Mezcla segmentos con distinta evolución de pista: Stroll (eliminado en Q2) queda por delante de Tsunoda porque su Q1 fue mejor. (3) Barras verticales con 20 nombres girados y etiquetas solapadas & Barras horizontales (como G16); nombrar la métrica «mejor vuelta de la sesión»; opción «por segmento» (Q1/Q2/Q3 por separado); marcar la línea de corte de Q1 y Q2 \\
G4 Lap chart & \textbf{Muy útil} & Etiqueta directa; parrilla en la vuelta 0; tabla con vueltas lideradas & Madeja en las ventanas de paradas (vueltas 10-17 y 30-38 en Bahrain); con 20 pilotos no se sigue a nadie salvo al que se resalta; no marca SC, abandonos ni paradas; datos desde 1950, pero en los años 50 solo el 56 \% de las posiciones por vuelta tienen dato, así que las líneas salen rotas & Resaltar por defecto el top 3 y atenuar al resto; marcar las paradas con puntos y el SC/VSC con bandas; ofrecer al lado un «race trace» (P1b), que responde mejor a la pregunta «quién ganó por estrategia» \\
G5 Stints & \textbf{Muy útil} & Claro y fiel a la convención de la TV; letra en cada tramo; el duro con contorno & Eje hasta 60 cuando la carrera tiene 57 vueltas; no distingue neumático nuevo o usado (el 20 \% de los stints de 2024 salen con más de 1 vuelta de uso, según \code{tyre\_\allowbreak{}age\_\allowbreak{}laps}); el primer stint de 1 vuelta (Hülkenberg) sale sin letra & \code{max =\allowbreak{} vueltas}; trama o transparencia para los usados; tooltip con la edad inicial \\
G6 Matriz & \textbf{Poco útil} (repetida) & Réplica fiel del TFG; accesible & Da la misma información que G5 en 20 × 57 celdas; obliga a desplazarse en horizontal; el recuadro de la vuelta de entrada no añade nada a los cortes de G5 & Fusionarla con G5 como su «tabla de datos», o reutilizar la forma de la matriz para algo nuevo: mapa de calor del \textbf{tiempo por vuelta respecto a la mediana} o de la \textbf{posición} (P3) \\
G7 Tiempos por vuelta & \textbf{Útil} para 2-4 pilotos, \textbf{poco útil} con 20 & Zoom doble (fiel al TFG); ocultar neutralizadas activado si hace falta; tabla completa & Con 20 líneas y 8 colores es ilegible (captura); las vueltas de boxes cortan las líneas; el eje de tiempos no es invertido (en F1 «más rápido» se suele dibujar arriba o se dibuja el delta) & Por defecto, solo el top 3 o los dos del mismo equipo; añadir la vista «delta con un piloto de referencia»; opción de corregir el combustible (P2) \\
G8 Violín & \textbf{Muy útil} & Filtro de vueltas cuidado (incluye la heurística del 150 \%); orden por mediana; tabla & Colores repetidos sin otra diferencia; texto diminuto; ticks con valores extraños (1:41.470, 1:39.698…) en lugar de valores redondos; la cola de VER hasta 1:32.6 (vuelta rápida con blandos nuevos) alarga mucho la escala & Ticks redondos (cada 0,5 s o 1 s); color de equipo; nombres horizontales más grandes; opción de dividir el violín por compuesto (mitad izquierda y derecha, \emph{split violin}) (P4) \\
G9 Media de paradas & \textbf{Poco útil tal como está} & Réplica del TFG; umbral de 60 s & La «duración» es el \textbf{tiempo en el pit lane} (mediana 2020s: 23,9 s), no la parada estática (unos 2-3 s), y el título no lo aclara; el eje en 0 aplana las diferencias; la media esconde las paradas malas & Renombrarlo a «Tiempo en el pit lane»; usar \emph{strip plot} (una marca por parada) con la mediana; dato que falta: la parada estática (DHL Fastest Pit Stop o documentos de la FIA) \\
G10 Telemetría & \textbf{Útil} (con un fallo de relato) & Tres paneles enlazados; selector A/B sin JS & Es la \textbf{vuelta rápida de clasificación} dentro de una pestaña de carrera y el título no lo dice; no hay \textbf{delta de tiempo acumulado} (lo más informativo en una comparación de telemetría); \code{rpm}, \code{gear} y \code{drs} están al 100 \% en gold pero no se usan; velocidad desde 0; el eje Y del freno sin sentido; «1,000 m» & Titular «Vuelta rápida de clasificación»; añadir el panel de delta (se integra con distancia y velocidad) y el de marchas; ajustar el eje de velocidad; ver P6 y P7 \\
G11 Trazado & \textbf{Útil} (atractivo) & Proporciones reales; escala continua & Solo el piloto A; no enseña la comparación que se acaba de elegir; la escala azul→naranja no marca dónde está la curva lenta & «Dominio por minisectores» A frente a B (P6) y el número de curva o la marcha en el trazado \\
G12 Puntos por temporada (piloto) & \textbf{Útil} & Título en morado; posición encima & Eje de categorías (se pierden los huecos); puntos no comparables entre eras; la barra no dice con qué equipo corrió & Eje temporal; alternar puntos, posición final y \% de puntos posibles; color por equipo (P13) \\
G13 Puntos frente a compañeros & \textbf{Poco útil} (repetido) & Textura accesible & Repite lo que ya dicen G14 y la tabla; al sumar a varios compañeros en una temporada (1994: Verstappen + Lehto + Herbert) se mezclan duelos distintos & Fusionar con G14 en un solo gráfico de barras divergentes por temporada y compañero (P12) \\
G14 H2H clasificación y carrera (×2) & \textbf{Útil} & Porcentajes claros & Dos gráficos casi iguales uno al lado del otro; el apilado al 100 \% gasta la mitad de la tinta en el «complemento» & Un único gráfico divergente: clasificación a la izquierda y carrera a la derecha, o 2 puntos por temporada \\
G15 Puntos por temporada (constructor) & \textbf{Poco útil} en equipos históricos & — & Ferrari 1950-57: el campeonato de constructores empezó en 1958, así que salen barras a 0 (el \code{points ?? 0} del código) con 3-7 victorias por temporada; puntos no comparables & Para antes de 1958, mostrar victorias o puntos de sus pilotos o un marcador «sin campeonato»; métrica normalizada \\
G16 Títulos & \textbf{Útil} & Horizontal, legible & Repite la columna «Títulos» de la tabla ordenable de al lado; el resumen no detecta el empate & Sustituirlo por algo que la tabla no dé: «títulos a lo largo del tiempo» (línea de tiempo por piloto, P16) \\
KPI temporada & \textbf{Poco útil} & Réplica fiel del TFG & El campeón sale 3 veces (tarjeta, pill en la tabla y resumen de G2); «victorias del campeón» es un dato pobre & Sustituirlas por KPIs de relato: margen del título, ronda de la decisión, nº de ganadores distintos, dominio del equipo (\% de victorias) \\
KPI piloto y constructor & \textbf{Útil} & Resumen rápido & Sin contexto (¿91 victorias es mucho?) & Añadir el percentil o el puesto histórico («2.º de todos los tiempos») o el porcentaje (\code{win\_\allowbreak{}rate\_\allowbreak{}pct} y \code{podium\_\allowbreak{}rate\_\allowbreak{}pct} existen en \code{agg\_\allowbreak{}driver\_\allowbreak{}career}) \\
Tabla de pasos por el pit lane & \textbf{Muy útil} y original & Distingue parada, SC, bandera roja y abandono & Poco visible, al final de la pestaña & Integrarla como marcas en G4 y G5 \\
\code{/\allowbreak{}quality} & \textbf{Útil} para el TFG y la transparencia & — & Solo tablas & Opcional: gráfico de barras de \% de coincidencia con la línea del umbral \\
\end{longtable}}

\subsection{Repeticiones y solapamientos (qué fusionar o eliminar)}

{\small\sloppy\setlength{\tabcolsep}{5pt}\renewcommand{\arraystretch}{1.15}
\begin{longtable}{@{}>{\raggedright\arraybackslash}p{0.192\dimexpr(\linewidth-30pt)\relax}>{\raggedright\arraybackslash}p{0.404\dimexpr(\linewidth-30pt)\relax}>{\raggedright\arraybackslash}p{0.404\dimexpr(\linewidth-30pt)\relax}@{}}
\toprule
\textbf{Solapamiento} & \textbf{Elementos} & \textbf{Propuesta} \\
\midrule\endhead
\bottomrule\endlastfoot
\textbf{Estrategia de neumáticos} & G5 stints + G6 matriz + tabla de stints & Dejar G5. La matriz pasa a ser su tabla alternativa o se reutiliza como mapa de calor del ritmo (P3) \\
\textbf{Ritmo de carrera} & G4 posiciones + G7 tiempos + G8 violín & Cada uno responde a algo distinto, pero G7 con 20 pilotos no responde a nada. Propuesta: pestaña «Carrera» con G4 y el race trace (P1b); pestaña «Ritmo» con el violín y G7 limitado a 2-4 pilotos o en modo delta \\
\textbf{Compañeros de equipo} & G13 + G14 ×2 + tabla por temporada (repetida 3 veces como tabla de cada figura) + tabla de detalle & Un único gráfico divergente (P12) + una tabla de detalle \\
\textbf{Campeón de la temporada} & 3 tarjetas + pill de la tabla + resumen de G2 + tarjeta en \code{/\allowbreak{}seasons} & Tarjetas de relato (ver KPI temporada) \\
\textbf{\% en clasificación} & G3 + columna \% de la tabla de clasificación & Aceptable (la tabla es la referencia), pero G3 debería aportar algo más (segmentos, cortes) \\
\textbf{Títulos} & G16 + columna «Títulos» del ranking & Sustituir G16 (P16) \\
\textbf{Inicio frente a temporada} & Top 5 de pilotos y constructores del inicio = tabla de la temporada actual & Aceptable como portada; mejor un «mini G2» (la lucha por el título en una línea) \\
\textbf{Puntos por temporada} & G12 (piloto) y G15 (constructor) con el mismo componente & Correcto (reutilización); arreglar los ejes y la normalización \\
\textbf{Salida y llegada} & Tabla de G4 (salida, llegada) y tabla de resultado (parrilla, posición) & Aceptable \\
\end{longtable}}

\subsection{Huecos de relato (preguntas típicas de un aficionado sin respuesta)}

\begin{enumerate}
\item \textbf{«¿Por cuánto ganó y cuándo se decidió la carrera?»} No hay diferencias en el tiempo. \code{gap\_\allowbreak{}to\_\allowbreak{}leader\_\allowbreak{}ms} está al 84-100 \% desde 2022, y el tiempo acumulado (Σ \code{lap\_\allowbreak{}time\_\allowbreak{}ms}) lo reconstruye exactamente desde 1996: en Bahrain 2024, vuelta 30, el cálculo coincide con la fuente al centésimo.
\item \textbf{«¿Hubo adelantamientos? ¿Fue una carrera entretenida?»} No hay índice de cambios de posición ni de cambios de líder (calculable: media de cambios de líder por carrera de 4,0 en los 50, 1,8 en los 70 y 2,7 en los 2020).
\item \textbf{«¿Qué pasa en este circuito?»} No hay \textbf{página de circuito}, aunque \code{dim\_\allowbreak{}race} tiene 78 circuitos con lat/lon, tipo, sentido, longitud y curvas. No se ven ganadores por circuito, conversión de la pole en victoria (Yas Marina 70,6 \%, Catalunya 69,4 \%…), probabilidad de SC ni récord de vuelta.
\item \textbf{«¿Quién es mejor: X o Y?»} No hay un comparador libre de pilotos. Solo se compara con los compañeros o, en telemetría, en una sesión.
\item \textbf{«¿Cuándo se decidió el título y quién podía ganarlo todavía?»} Hay 76 carreras marcadas como \code{drivers\_\allowbreak{}championship\_\allowbreak{}decider}, pero no se usan. Tampoco se calcula la eliminación matemática.
\item \textbf{Fiabilidad y abandonos por época.} \code{reason\_\allowbreak{}retired} tiene datos de todas las décadas. Los clasificados pasan del 41 \% en los 80 al 87 \% en los 2020, y la parte mecánica de los abandonos baja del 86 \% al 55 \%.
\item \textbf{Edades y carreras deportivas.} Ganador más joven (Verstappen, 18,6 años) y edad media de la parrilla (35,1 en los 50 frente a 28,4 en los 2020). \code{date\_\allowbreak{}of\_\allowbreak{}birth} está al 100 \%.
\item \textbf{Nacionalidades y «victorias en casa».} Hay 95 victorias en casa de 1 167.
\item \textbf{Equipos y motores.} No se ve el motor por equipo y temporada, aunque \code{engine\_\allowbreak{}manufacturer\_\allowbreak{}id} está en los resultados. Tampoco hay una línea de tiempo de las alineaciones.
\item \textbf{Sprint.} Hay resultado de sprint, pero no vueltas (ver nota 03).
\item \textbf{Datos que no tenemos.} Meteorología, mensajes de dirección de carrera (solo hay banderas por vuelta), parada estática y adelantamientos oficiales.
\end{enumerate}

\medskip\noindent\textcolor{muted!40}{\rule{\linewidth}{0.3pt}}\medskip

\section{Propuestas nuevas}

\subsection{Qué se puede dibujar y desde cuándo (comprobado en \code{data/\allowbreak{}gold/\allowbreak{}f1.\allowbreak{}duckdb})}

{\small\sloppy\setlength{\tabcolsep}{5pt}\renewcommand{\arraystretch}{1.15}
\begin{longtable}{@{}>{\raggedright\arraybackslash}p{0.253\dimexpr(\linewidth-30pt)\relax}>{\raggedright\arraybackslash}p{0.373\dimexpr(\linewidth-30pt)\relax}>{\raggedright\arraybackslash}p{0.373\dimexpr(\linewidth-30pt)\relax}@{}}
\toprule
\textbf{Dato} & \textbf{Tabla y columna} & \textbf{Cobertura} \\
\midrule\endhead
\bottomrule\endlastfoot
Posición por vuelta & \code{fact\_\allowbreak{}laptimes.\allowbreak{}position} & Las 1 164 carreras disputadas desde 1950 (56 \% de las filas con dato en los 50, 95 \% en los 60, 100 \% desde los 70) \\
Tiempo por vuelta & \code{fact\_\allowbreak{}laptimes.\allowbreak{}lap\_\allowbreak{}time\_\allowbreak{}ms} & 0 \% antes de 1990, 96,5 \% en los 90, \textasciitilde{}100 \% desde 2000 (desde 1996 en la práctica) \\
Gap al líder & \code{gap\_\allowbreak{}to\_\allowbreak{}leader\_\allowbreak{}ms} & 60-87 \% por temporada desde 1996, 100 \% en 2025-26; \textbf{reconstruible} exactamente con Σ \code{lap\_\allowbreak{}time\_\allowbreak{}ms} (comprobado en Bahrain 2024, vuelta 30: Pérez +15,68 s, Sainz +18,20 s en ambos cálculos) \\
Compuesto y edad del neumático & \code{tyre\_\allowbreak{}compound}, \code{tyre\_\allowbreak{}age\_\allowbreak{}laps} & Desde 2011 (100 \%); compuesto Pirelli C1-C5 desde 2019 \\
Sectores y speed trap & \code{sector\_\allowbreak{}*\_\allowbreak{}ms}, \code{speed\_\allowbreak{}trap\_\allowbreak{}kmh} & Desde 2018 (98 \%) \\
Banderas por vuelta & \code{is\_\allowbreak{}safety\_\allowbreak{}car}, \code{is\_\allowbreak{}yellow\_\allowbreak{}flag}, \code{is\_\allowbreak{}virtual\_\allowbreak{}safety\_\allowbreak{}car}, \code{is\_\allowbreak{}red\_\allowbreak{}flag} & SC desde los 90 (1 922 vueltas), VSC desde 2015, bandera roja desde 2000 \\
Paradas & \code{fact\_\allowbreak{}pit\_\allowbreak{}stops.\allowbreak{}time\_\allowbreak{}ms} & Desde 1994 (612 carreras); es el tiempo \textbf{en el pit lane} (mediana de 30,0 s en los 90 y 23,9 s en los 2020) \\
Pasos por el pit lane & \code{fact\_\allowbreak{}pit\_\allowbreak{}lane\_\allowbreak{}passes.\allowbreak{}pass\_\allowbreak{}type} & pit\_stop 22 485, unclassified 1 118, retirement 351, penalty\_or\_other 141, safety\_car 114, red\_flag 62 \\
Telemetría de clasificación & \code{fact\_\allowbreak{}quali\_\allowbreak{}telemetry} & 2024-2026 (63 carreras); \textasciitilde{}450 puntos por vuelta; velocidad, rpm, marcha, acelerador, freno, DRS y x/y al 100 \% \\
Clasificación del campeonato por ronda & \code{fact\_\allowbreak{}driver\_\allowbreak{}standing}, \code{fact\_\allowbreak{}constructor\_\allowbreak{}standing} & Todas las rondas desde 1950 (1950: 108/108 filas con posición) \\
Resultado & \code{fact\_\allowbreak{}race\_\allowbreak{}result} & Parrilla \textasciitilde{}90 \%; \code{reason\_\allowbreak{}retired} con 198 valores distintos; \code{positions\_\allowbreak{}gained}; \code{is\_\allowbreak{}grand\_\allowbreak{}slam} (71); piloto del día (2016+, 228) \\
Q1/Q2/Q3 & \code{fact\_\allowbreak{}qualifying\_\allowbreak{}result} & Desde 2006 (40 \% de los 2000, \textasciitilde{}99 \% desde 2010); mejor tiempo desde 1950 (96 \%) \\
Carrera & \code{dim\_\allowbreak{}race} & 78 circuitos con lat/lon, tipo (ROAD/RACE/STREET), sentido, longitud, curvas; \code{drivers\_\allowbreak{}championship\_\allowbreak{}decider} (76 carreras) \\
Piloto & \code{dim\_\allowbreak{}driver} & \code{date\_\allowbreak{}of\_\allowbreak{}birth} 100 \% (917 pilotos), nacionalidad alpha2 \\
Compañeros & \code{agg\_\allowbreak{}teammate\_\allowbreak{}h2h} & 13 530 filas; grafo de 4 componentes conexas, con 810 pilotos en la mayor \\
\end{longtable}}

\textbf{Datos que faltan} para algunas propuestas:

\begin{itemize}
\item \textbf{Colores de equipo por temporada}: hace falta un seed. La vista \code{stg\_\allowbreak{}fastf1\_\allowbreak{}\_\allowbreak{}results.\allowbreak{}team\_\allowbreak{}color} apunta a un bronze que no está disponible y solo cubriría 2018+.
\item \textbf{Clasificación de \code{reason\_\allowbreak{}retired}} en mecánico, accidente u otro: seed de unas 200 filas.
\item \textbf{Parada estática} (solo tenemos el tiempo en el pit lane).
\item \textbf{Meteorología}.
\item \textbf{Mensajes de dirección de carrera} (solo hay banderas por vuelta).
\item \textbf{Vueltas de sprint} (ver nota 03).
\item \textbf{Adelantamientos oficiales}: los cambios de posición son un proxy.
\end{itemize}

Esfuerzo: \textbf{S} = menos de 1 día (reutiliza un endpoint o componente existente); \textbf{M} = 1-3 días (endpoint o modelo dbt nuevo); \textbf{L} = más de 3 días.

\subsection{Carrera (pestañas de \code{/\allowbreak{}races/\allowbreak{}[id]})}

\textbf{P1b. Race trace: diferencia con el líder vuelta a vuelta.} \emph{Esfuerzo S-M.}

\begin{itemize}
\item Pregunta: ¿por cuánto iba ganando, cuándo se abrió o cerró la carrera y qué hizo el SC con las diferencias?
\item Boceto:
\begin{itemize}
\item Eje X: vuelta. Eje Y: segundos detrás del líder, invertido, con el líder en 0 arriba.
\item Una línea por piloto; las vueltas bajo SC/VSC como bandas grises de fondo y las paradas como puntos huecos.
\item Etiqueta directa al final. Por defecto se resalta el top 5 y el resto va en gris al 25 \%.
\item Variante: «delta con un piloto elegido» (la referencia horizontal es ese piloto).
\end{itemize}
\item Datos: \code{fact\_\allowbreak{}laptimes} (Σ \code{lap\_\allowbreak{}time\_\allowbreak{}ms} por piloto y vuelta; en cada vuelta, referencia = mínimo acumulado), \code{is\_\allowbreak{}safety\_\allowbreak{}car}/\code{is\_\allowbreak{}virtual\_\allowbreak{}safety\_\allowbreak{}car}, \code{is\_\allowbreak{}pit\_\allowbreak{}in\_\allowbreak{}lap}. Desde 1996.
\item Dónde: pestaña «Vuelta a vuelta», junto a G4, con un conmutador «posiciones | diferencias».
\item Accesibilidad: resumen generado («X abrió N s en M vueltas; el SC de la vuelta K dejó la ventaja en 1 s») y tabla vuelta × piloto.
\end{itemize}

\textbf{P2. Degradación por stint.} \emph{Esfuerzo M.}

\begin{itemize}
\item Pregunta: ¿quién cuidó mejor los neumáticos? ¿Cuánto perdía por vuelta cada compuesto?
\item Boceto:
\begin{itemize}
\item (a) \emph{Dot plot}: una fila por piloto; cada stint es un punto en X = pendiente en ms/vuelta, con color y letra del compuesto.
\item (b) Al pulsar un piloto, dispersión de tiempo corregido frente a la edad del neumático, con su recta de regresión por stint.
\item Corrección de combustible: restar la tendencia común de la carrera (mediana por vuelta de todos los pilotos) o un coeficiente fijo documentado como estimación.
\end{itemize}
\item Datos: \code{regr\_\allowbreak{}slope(\allowbreak{}lap\_\allowbreak{}time\_\allowbreak{}ms,\allowbreak{} tyre\_\allowbreak{}age\_\allowbreak{}laps)} por (piloto, stint) sobre vueltas limpias (sin vuelta 1, boxes ni SC/VSC; mínimo 8 vueltas). Probado en Bahrain 2024: pendientes de 6 a 123 ms/vuelta (Leclerc: 84 ms/vuelta con blando y 6 con el primer duro). La dispersión muestra que la corrección de combustible es imprescindible antes de publicarlo. Desde 2011.
\item Dónde: pestaña «Neumáticos» o «Estrategia», \textbf{en lugar de la matriz G6}.
\item Accesibilidad: tabla piloto/stint/compuesto/vueltas/pendiente; resumen «el que menos degradó con duro fue X».
\end{itemize}

\textbf{P3. Mapa de calor del ritmo (reutiliza la forma de la matriz).} \emph{Esfuerzo S.}

\begin{itemize}
\item Pregunta: ¿en qué fase de la carrera fue rápido cada piloto?
\item Boceto:
\begin{itemize}
\item Filas = pilotos (orden de llegada); columnas = vueltas.
\item Color divergente = diferencia con la mediana de la vuelta, con los rápidos en azul y los lentos en naranja (apto para daltonismo).
\item Vueltas neutralizadas rayadas, paradas recuadradas y la letra del compuesto opcional.
\end{itemize}
\item Datos: los mismos \code{laps} que usa hoy G6.
\item Dónde: sustituye a G6.
\item Accesibilidad: sigue siendo una \code{<table>} con el valor en texto \code{sr-\allowbreak{}only}, igual que la matriz actual.
\end{itemize}

\textbf{P4. Violín partido por compuesto.} \emph{Esfuerzo S.}

\begin{itemize}
\item Pregunta: ¿el ritmo de X era mejor con medio o con duro?
\item Boceto: mitad izquierda del violín con el compuesto 1 y mitad derecha con el compuesto 2, en colores Pirelli con contorno; ticks redondos.
\item Datos: los de G8 más \code{tyre\_\allowbreak{}compound}.
\item Dónde: opción de G8.
\item Accesibilidad: la tabla gana una columna por compuesto.
\end{itemize}

\textbf{P5. Detector de undercut y overcut.} \emph{Esfuerzo M.} No lo he visto como gráfico automático en otras webs de aficionados.

\begin{itemize}
\item Pregunta: ¿quién ganó o perdió posiciones por la estrategia de paradas?
\item Boceto:
\begin{itemize}
\item Lista de «duelos de estrategia»; cada duelo es un mini \emph{slope chart} con dos puntos (antes y después) y dos líneas (A y B).
\item Etiqueta: «Norris paró 2 vueltas antes que Leclerc → undercut logrado (+1 posición)».
\item Verde u óxido según el resultado, con icono \si{}/\no{} para no depender del color.
\end{itemize}
\item Datos:
\begin{itemize}
\item Heurística: A y B van consecutivos en pista la vuelta anterior a la parada de A; B para entre 1 y 3 vueltas después; se comparan las posiciones en la vuelta \code{lb+2}.
\item Resultado en 2024: \textbf{204 intentos, 49 undercuts logrados}.
\item Mejora: pedir que la diferencia previa sea menor de 3 s (con \code{gap\_\allowbreak{}to\_\allowbreak{}leader\_\allowbreak{}ms} o el acumulado) y descartar las paradas bajo SC.
\item Desde 1996 (las paradas empiezan en 1994).
\end{itemize}
\item Dónde: pestaña «Boxes» o «Estrategia» (le da el relato que ahora le falta).
\item Accesibilidad: es una lista de frases; el gráfico solo ilustra.
\end{itemize}

\textbf{P6. Delta de tiempo y dominio por minisectores (telemetría).} \emph{Esfuerzo M.}

\begin{itemize}
\item Pregunta: ¿en qué curvas gana tiempo A sobre B?
\item Boceto:
\begin{itemize}
\item (a) Panel nuevo encima de la velocidad: delta acumulado (s) frente a la distancia, con el cero en el centro; por encima, A es más lento.
\item (b) Trazado dividido en unos 25 minisectores, cada uno con el color del piloto más rápido en él y patrón o letra A/B para no depender del color. Es el gráfico más conocido de FastF1, pero no hay una web en español que lo ofrezca navegable.
\end{itemize}
\item Datos: \code{fact\_\allowbreak{}quali\_\allowbreak{}telemetry} (distancia, velocidad, x, y). El tiempo acumulado se obtiene con t = Σ Δd/v y se interpola a una rejilla común de distancia. 2024+.
\item Dónde: telemetría, que debería pasar a la pestaña Clasificación (sección 6).
\item Accesibilidad: tabla por minisector (distancia, piloto más rápido, diferencia).
\end{itemize}

\textbf{P7. Mapa de marchas y DRS en el trazado.} \emph{Esfuerzo S.}

\begin{itemize}
\item Pregunta: ¿cómo se conduce este circuito?
\item Boceto: trazado coloreado por marcha (paleta ordinal de 8 pasos con el número en las zonas largas), tramos con DRS abierto en trazo grueso y RPM en el tooltip.
\item Datos: \code{gear}, \code{drs}, \code{rpm} (100 \%, sin usar hoy). 2024+.
\item Dónde: telemetría y página de circuito (P24).
\end{itemize}

\textbf{P8. «Ghost race» animada.} \emph{Esfuerzo L.} Inédita en formato accesible.

\begin{itemize}
\item Pregunta: ¿cómo se vería la carrera desde arriba?
\item Boceto:
\begin{itemize}
\item Un punto por coche (código y color de equipo) sobre la silueta del circuito, o sobre una recta horizontal que representa una vuelta cuando no hay trazado.
\item Posición = fracción de vuelta según el tiempo acumulado interpolado.
\item Controles: play/pausa, velocidad ×10/×30 y deslizador de vuelta. Bandas SC.
\end{itemize}
\item Datos: Σ \code{lap\_\allowbreak{}time\_\allowbreak{}ms} (desde 1996); silueta de \code{fact\_\allowbreak{}quali\_\allowbreak{}telemetry} x/y (2024+) o, antes, una línea neutra.
\item Accesibilidad:
\begin{itemize}
\item Con \code{prefers-\allowbreak{}reduced-\allowbreak{}motion} se muestra en su lugar una serie estática de mini-gráficos (vueltas 1, 10, 20…).
\item El deslizador es un \code{<input type=\allowbreak{}"range">} con \code{aria-\allowbreak{}valuetext} («Vuelta 23: 1.º VER, 2.º PER a 4,1 s…»).
\item Nunca arranca sola.
\end{itemize}
\end{itemize}

\textbf{P9. Parrilla → vuelta 1 → llegada (\emph{slope chart}).} \emph{Esfuerzo S.}

\begin{itemize}
\item Pregunta: ¿quién ganó posiciones en la salida y quién en la carrera?
\item Boceto: 3 columnas (parrilla, fin de la vuelta 1 y llegada); una línea por piloto, verde si gana y óxido si pierde, con grosor o flecha como segunda señal; los abandonos terminan en una columna «Ab.».
\item Datos: \code{grid\_\allowbreak{}position}, \code{fact\_\allowbreak{}laptimes.\allowbreak{}position} de la vuelta 1 y \code{position\_\allowbreak{}number}. \textbf{Desde 1950.}
\item Dónde: pestaña «Resultado», encima de la tabla.
\item Extra histórico: ranking de «mejores salidas» (mínimo de 100 carreras: Jochen Mass +1,56 posiciones de media en la vuelta 1, Brundle +1,52, Stroll +1,34).
\end{itemize}

\subsection{Temporada (\code{/\allowbreak{}seasons/\allowbreak{}[year]})}

\textbf{P1. Diferencia con el líder del campeonato y eliminación matemática.} \emph{Esfuerzo S-M.}

\begin{itemize}
\item Pregunta: ¿cuándo se decidió el título y quién seguía con opciones?
\item Boceto:
\begin{itemize}
\item Eje X: ronda. Eje Y: puntos detrás del líder, invertido.
\item Banda sombreada con los «puntos que quedan en juego» (restantes × máximo por carrera; en gold, el máximo es 9 en 1950 y 1990, 25 en 2010 y 2025 y 26 en 2019 por la vuelta rápida, más el sprint).
\item Cuando la línea de un piloto sale de la banda, un \no{} indica que queda eliminado. Una línea vertical marca la carrera \code{drivers\_\allowbreak{}championship\_\allowbreak{}decider}.
\item Etiquetas directas con separación real (hay que resolver el solape de G2).
\end{itemize}
\item Datos: \code{fact\_\allowbreak{}driver\_\allowbreak{}standing} (todas las rondas desde 1950), \code{fact\_\allowbreak{}race\_\allowbreak{}result.\allowbreak{}points} para el máximo por carrera y \code{dim\_\allowbreak{}race.\allowbreak{}drivers\_\allowbreak{}championship\_\allowbreak{}decider}. Aviso: las reglas de «mejores N resultados» anteriores a 1991 hacen que la eliminación sea aproximada y hay que indicarlo.
\item Dónde: modo alternativo de G2 (conmutador «puntos | diferencia | posición»).
\item Accesibilidad: resumen «El título se decidió en la ronda N (GP de X); Y quedó eliminado en la ronda M».
\end{itemize}

\textbf{P10. Bump chart de posiciones en el campeonato.} \emph{Esfuerzo S.}

\begin{itemize}
\item Pregunta: ¿cómo cambió el orden del campeonato ronda a ronda?
\item Boceto: como G4, pero por ronda y con \code{position\_\allowbreak{}number} de la clasificación del campeonato.
\item Datos: \code{fact\_\allowbreak{}driver\_\allowbreak{}standing.\allowbreak{}position\_\allowbreak{}number}.
\item Dónde: tercera opción del conmutador de G2. Reutiliza \code{Lap\allowbreak{}Chart}.
\end{itemize}

\textbf{P11. KPIs de relato de la temporada.} \emph{Esfuerzo S.} Sustituyen a las tarjetas actuales.

\begin{itemize}
\item Contenido: margen final del título, ronda en que se decidió, nº de ganadores distintos, \% de victorias del mejor equipo y, en temporadas en curso, «carreras restantes / puntos en juego».
\item Datos: standings, \code{is\_\allowbreak{}win} y \code{drivers\_\allowbreak{}championship\_\allowbreak{}decider}.
\end{itemize}

\subsection{Piloto y constructor}

\textbf{P12. H2H divergente con los compañeros (fusiona G13 y G14).} \emph{Esfuerzo S.}

\begin{itemize}
\item Pregunta: ¿batió a sus compañeros, y a cuál?
\item Boceto:
\begin{itemize}
\item Una fila por temporada y compañero; barra divergente desde el 50 \%: a la derecha, \% de clasificaciones ganadas; a la izquierda, perdidas.
\item Un segundo marcador (rombo) con el \% en carrera.
\item Etiqueta con el nombre del compañero y color del equipo.
\item Se ve de un vistazo, por ejemplo, el 2010-2012 de Schumacher frente a Rosberg.
\end{itemize}
\item Datos: \code{agg\_\allowbreak{}teammate\_\allowbreak{}h2h} (ya en \code{/\allowbreak{}drivers/\allowbreak{}\{id\}/\allowbreak{}teammates}).
\item Accesibilidad: una única tabla (hoy la misma tabla aparece 3 veces).
\end{itemize}

\textbf{P13. «ADN» o carrera deportiva del piloto.} \emph{Esfuerzo M.}

\begin{itemize}
\item Pregunta: ¿qué tipo de piloto es: poleman, de domingo, fiable, remontador?
\item Boceto:
\begin{itemize}
\item (a) Franja temporal (\emph{swimlane}) sobre un eje de años continuo: las temporadas coloreadas por equipo y un punto con la posición final (eje invertido). Se ven los huecos, como la retirada de Schumacher de 2007 a 2009.
\item (b) «Huella» con 6 barras de percentil frente a los pilotos de su época (no radar, que distorsiona las áreas): \% de victorias, \% de podios, \% de poles, H2H en clasificación, posiciones ganadas de media y \% de abandonos. Cada barra lleva la cifra real.
\end{itemize}
\item Datos: \code{agg\_\allowbreak{}driver\_\allowbreak{}career} (\code{win\_\allowbreak{}rate\_\allowbreak{}pct}, \code{podium\_\allowbreak{}rate\_\allowbreak{}pct}, \code{pole\_\allowbreak{}positions}, \code{race\_\allowbreak{}starts}), \code{agg\_\allowbreak{}teammate\_\allowbreak{}h2h}, \code{fact\_\allowbreak{}race\_\allowbreak{}result.\allowbreak{}positions\_\allowbreak{}gained} y \code{reason\_\allowbreak{}retired}. Para los percentiles, pilotos con más de 30 salidas en la misma década.
\item Accesibilidad: tabla métrica/valor/percentil.
\end{itemize}

\textbf{P14. Comparación de eras normalizada.} \emph{Esfuerzo S.}

\begin{itemize}
\item Pregunta: ¿quién habría sumado más con el sistema de puntos actual?
\item Boceto: en \code{/\allowbreak{}records} y en la ficha, conmutador «puntos oficiales | sistema 2010 (25-18-15-12-10-8-6-4-2-1) | puntos por carrera».
\item Datos: se recalcula desde \code{fact\_\allowbreak{}race\_\allowbreak{}result.\allowbreak{}position\_\allowbreak{}number}. Comprobado (mínimo de 50 carreras), en puntos 2010 por carrera: Fangio 15,05, Verstappen 13,86, Hamilton 13,85, M. Schumacher 12,63, Prost 12,26, Clark 11,34, Senna 11,61 y Stewart 11,09.
\item Dónde: \code{/\allowbreak{}records} (nueva métrica ordenable) y G12 (eje alternativo).
\item Accesibilidad: es una columna más de la tabla.
\end{itemize}

\textbf{P15. Red de compañeros y «grados de separación».} \emph{Esfuerzo M.} Inédito.

\begin{itemize}
\item Pregunta: ¿cómo se conecta Fangio con Antonelli?
\item Boceto:
\begin{itemize}
\item (a) Buscador de dos pilotos → camino mínimo dibujado como cadena horizontal de «fichas» (piloto → equipo y año → piloto…).
\item (b) Grafo de fuerzas opcional del vecindario a 2 saltos (no los 810 nodos a la vez).
\end{itemize}
\item Datos: aristas de \code{agg\_\allowbreak{}teammate\_\allowbreak{}h2h}. BFS en Python sobre gold: \textbf{Fangio → Jo Bonnier → Rolf Stommelen → Riccardo Patrese → Michael Schumacher → Nico Rosberg → Lewis Hamilton → George Russell → Kimi Antonelli} (8 saltos). 4 componentes conexas; la mayor tiene 810 pilotos. Con más de 900 nodos, el BFS se calcula en milisegundos en la API (Python) o se precalcula. La consulta recursiva en SQL explotaba combinatoriamente y no conviene.
\item Accesibilidad: el camino es una lista ordenada de texto; el grafo es decorativo, con alternativa textual.
\end{itemize}

\textbf{P16. Línea de tiempo de títulos (sustituye a G16).} \emph{Esfuerzo S.}

\begin{itemize}
\item Pregunta: ¿quién dominó cada época?
\item Boceto:
\begin{itemize}
\item Eje X: temporada 1950-2026. Una fila por campeón, ordenada por su primer título.
\item Un cuadrado por título con el color del equipo y una línea fina que une su primer y último título.
\item Filtro por rango (el que ya existe).
\end{itemize}
\item Datos: \code{agg\_\allowbreak{}driver\_\allowbreak{}season.\allowbreak{}championship\_\allowbreak{}won} (o \code{/\allowbreak{}seasons}).
\item Accesibilidad: tabla campeón → años.
\end{itemize}

\textbf{P17. Alineación y motor del constructor.} \emph{Esfuerzo S.}

\begin{itemize}
\item Pregunta: ¿quién pilotó para el equipo y con qué motor?
\item Boceto: \emph{swimlane} por temporada con una fila por piloto (barra por temporada); banda superior coloreada por motor.
\item Datos: \code{/\allowbreak{}constructors/\allowbreak{}\{id\}/\allowbreak{}seasons} (ya trae los pilotos) y \code{fact\_\allowbreak{}race\_\allowbreak{}result.\allowbreak{}engine\_\allowbreak{}manufacturer\_\allowbreak{}id} (22 fabricantes con victorias).
\item Arreglo relacionado: en G15, antes de 1958 no dibujar barras a 0, sino un marcador «sin campeonato de constructores» o las victorias.
\end{itemize}

\subsection{Historia (vistas agregadas nuevas)}

\textbf{P18. Streamgraph de victorias por constructor, 1950-2026.} \emph{Esfuerzo S-M.}

\begin{itemize}
\item Pregunta: ¿cómo han cambiado las eras de dominio?
\item Boceto:
\begin{itemize}
\item Eje X: temporada. El grosor de cada capa es el \% de victorias de la temporada (normalizado, para que las temporadas de 7 y de 24 carreras pesen igual).
\item 10 equipos con más victorias en color y el resto en «otros» gris; etiquetas dentro de las capas; conmutador «constructores | motores».
\item Línea superpuesta con el «índice de dominio» (\% de victorias del mejor equipo de cada año).
\end{itemize}
\item Datos: \code{fact\_\allowbreak{}race\_\allowbreak{}result.\allowbreak{}is\_\allowbreak{}win} × \code{constructor\_\allowbreak{}id}/\code{engine\_\allowbreak{}manufacturer\_\allowbreak{}id} (38 constructores y 22 motores con victorias; p. ej., 2025: McLaren 14, Red Bull 8, Mercedes 2).
\item Accesibilidad: tabla temporada × equipo y resumen por eras.
\item Nota: el streamgraph centrado es difícil de leer con precisión; la alternativa es un área apilada al 100 \%.
\end{itemize}

\textbf{P19. Calendario mundial: país × temporada.} \emph{Esfuerzo S.}

\begin{itemize}
\item Pregunta: ¿cuándo entró y salió la F1 de cada país?
\item Boceto:
\begin{itemize}
\item Rejilla con filas = países, agrupados por continente y ordenados por primer GP, y columnas = temporadas.
\item Celda rellena con el nº de GP de ese año (1, 2 o 3) y el circuito en el tooltip.
\item Parecido a un gráfico de contribuciones de GitHub.
\end{itemize}
\item Datos: \code{dim\_\allowbreak{}race} (\code{circuit\_\allowbreak{}country}/\code{grand\_\allowbreak{}prix\_\allowbreak{}country}, \code{season}, \code{circuit\_\allowbreak{}name}). El continente habría que añadirlo a gold: \code{stg\_\allowbreak{}f1db\_\allowbreak{}\_\allowbreak{}countries.\allowbreak{}continent\_\allowbreak{}id} existe, pero la vista depende del bronze.
\item Dónde: inicio, como segunda vista del mapa («mapa | calendario»). Responde a lo que el mapa no muestra: el tiempo.
\item Accesibilidad: es una tabla país × década con recuento.
\end{itemize}

\textbf{P20. Fiabilidad a lo largo de la historia.} \emph{Esfuerzo M} (por el seed).

\begin{itemize}
\item Pregunta: ¿son los coches más fiables? ¿Se abandona más por accidentes?
\item Boceto: área apilada al 100 \% por temporada con clasificados, abandono mecánico y abandono por accidente u otro, con anotaciones. Solo se clasificaba el 41 \% de los participantes en los 80 y el 87 \% en los 2020.
\item Datos: \code{reason\_\allowbreak{}retired} (198 valores) → seed \code{retirement\_\allowbreak{}categories.\allowbreak{}csv}. Con una clasificación provisional, la parte mecánica de los abandonos baja del 86 \% (años 60) al 55 \% (2020s) y la de accidentes sube del 12 \% al 44 \%.
\item Dónde: Historia. Versión por equipo en la ficha del constructor.
\end{itemize}

\textbf{P21. Mapa de calor de posiciones de llegada.} \emph{Esfuerzo S.}

\begin{itemize}
\item Pregunta: ¿cómo terminaba de verdad un piloto: siempre en puntos o todo o nada?
\item Boceto:
\begin{itemize}
\item En la ficha: filas = temporadas; columnas = posición de llegada 1…20 + «Ab.».
\item Intensidad secuencial = nº de veces; la columna 1 lleva el borde morado.
\item Versión global en Historia: filas = décadas.
\end{itemize}
\item Datos: \code{fact\_\allowbreak{}race\_\allowbreak{}result.\allowbreak{}position\_\allowbreak{}number} y \code{reason\_\allowbreak{}retired}. Desde 1950.
\item Accesibilidad: es una \code{<table>} con los recuentos en texto.
\end{itemize}

\textbf{P22. Edad y relevo generacional.} \emph{Esfuerzo S.}

\begin{itemize}
\item Pregunta: ¿los pilotos son cada vez más jóvenes?
\item Boceto: \emph{beeswarm} por temporada con la edad de cada ganador el día de la victoria; línea con la edad media de la parrilla (35,1 años en los 50 → 28,4 en los 2020); anotaciones para el más joven (Verstappen, 18,62 años en 2016; le sigue Antonelli con 19,55 en 2026) y el mayor.
\item Datos: \code{dim\_\allowbreak{}driver.\allowbreak{}date\_\allowbreak{}of\_\allowbreak{}birth} (100 \%) y \code{dim\_\allowbreak{}race.\allowbreak{}race\_\allowbreak{}date}.
\end{itemize}

\textbf{P23. Margen de victoria e «índice de emoción».} \emph{Esfuerzo M.}

\begin{itemize}
\item Pregunta: ¿las carreras son hoy más ajustadas? ¿Cuáles fueron las más disputadas?
\item Boceto:
\begin{itemize}
\item (a) \emph{Strip plot} del margen del ganador por carrera (escala logarítmica) con la mediana móvil: 27,6 s en los 50, 11,2 s en los 90 y 6,2 s en los 2020.
\item (b) Tabla «carreras más disputadas» con un índice compuesto de cambios de líder, cambios de posición en pista (sin la vuelta 1 ni vueltas de boxes) y margen final.
\item El índice es un \textbf{proxy} y hay que explicarlo: sin datos oficiales de adelantamientos, un cambio de posición no es siempre un adelantamiento.
\end{itemize}
\item Datos: \code{gap\_\allowbreak{}ms} de P2 y posición por vuelta. Cambios de líder por carrera: media de 4,0 en los 50, 1,8 en los 70, 3,6 en los 2000 y 2,7 en los 2020.
\item Accesibilidad: el ranking es una tabla.
\end{itemize}

\subsection{Entidades nuevas}

\textbf{P24. Página de circuito.} \emph{Esfuerzo M.} Es el hueco más grande.

\begin{itemize}
\item Pregunta: ¿qué pasa en este circuito?
\item Contenido:
\begin{itemize}
\item Ficha: tipo, sentido, longitud, curvas y temporadas.
\item Trazado de P7 si hay telemetría.
\item Ganadores por año (tira de cuadrados con el color del equipo) y conversión de pole en victoria (mínimo de 15 carreras: Yas Marina 70,6 \%, Catalunya 69,4 \%, Marina Bay 68,8 \%, Shanghái 63,2 \%).
\item \% de vueltas bajo SC desde 2000 (Marina Bay 9,8 \%, Baku 9,1 \%, Melbourne 8,5 \%, Interlagos 7,8 \%).
\item Récord de vuelta en carrera por trazado: hay que agrupar por \code{course\_\allowbreak{}length}, porque los trazados cambian.
\end{itemize}
\item Datos: \code{dim\_\allowbreak{}race}, \code{fact\_\allowbreak{}race\_\allowbreak{}result} y \code{fact\_\allowbreak{}laptimes}. Endpoint nuevo \code{/\allowbreak{}circuits/\allowbreak{}\{id\}}; la ruta \code{/\allowbreak{}circuits} ya existe para el mapa.
\item Dónde: los nombres de circuito de las tablas y del mapa pasan a ser enlaces.
\end{itemize}

\textbf{P25. Comparador libre de dos pilotos.} \emph{Esfuerzo M.}

\begin{itemize}
\item Pregunta: ¿quién es mejor, X o Y?
\item Boceto:
\begin{itemize}
\item (a) Victorias acumuladas frente al nº de carrera, o frente a la edad (dos líneas con etiqueta directa).
\item (b) KPIs enfrentados en espejo (barras a izquierda y derecha).
\item (c) Si coincidieron, H2H en carreras compartidas (quién terminó delante).
\item (d) Normalización de P14.
\end{itemize}
\item Datos: \code{fact\_\allowbreak{}race\_\allowbreak{}result}, \code{agg\_\allowbreak{}driver\_\allowbreak{}career} y \code{dim\_\allowbreak{}driver}.
\item Dónde: \code{/\allowbreak{}compare?a=\allowbreak{}\&b=\allowbreak{}}, enlazado desde cada ficha («comparar con…»).
\item Accesibilidad: formulario GET sin JS (como la telemetría) y una tabla comparativa.
\end{itemize}

\textbf{P26. Victorias en casa.} \emph{Esfuerzo S.}

\begin{itemize}
\item Pregunta: ¿rinden más los pilotos en su país?
\item Boceto: barras por nacionalidad con victorias en casa y fuera (95 de 1 167 victorias fueron en casa, comparando \code{nationality\_\allowbreak{}alpha2} con \code{circuit\_\allowbreak{}alpha2}).
\item Dónde: Historia o la página de circuito.
\end{itemize}

\subsection{Prioridad sugerida}

{\small\sloppy\setlength{\tabcolsep}{5pt}\renewcommand{\arraystretch}{1.15}
\begin{longtable}{@{}>{\raggedright\arraybackslash}p{0.170\dimexpr(\linewidth-30pt)\relax}>{\raggedright\arraybackslash}p{0.415\dimexpr(\linewidth-30pt)\relax}>{\raggedright\arraybackslash}p{0.415\dimexpr(\linewidth-30pt)\relax}@{}}
\toprule
\textbf{Prioridad} & \textbf{Propuestas} & \textbf{Motivo} \\
\midrule\endhead
\bottomrule\endlastfoot
1 (arreglos baratos) & Ejes temporales en G12, G13 y G15; ticks y texto del violín; formato numérico local; G3 horizontal y con la métrica bien nombrada; seed de colores de equipo; etiquetas de G2; títulos de G9 («tiempo en el pit lane») y G10 («vuelta de clasificación»); barras a 0 de G15; empates en los resúmenes & Errores de lectura con impacto directo y poco coste \\
2 (relato de carrera) & P1b race trace, P3 en lugar de la matriz, P6 delta y minisectores, P9 parrilla→llegada, P12 H2H divergente & Responden a las preguntas más frecuentes y usan datos que ya sirve la API \\
3 (entidades) & P24 circuito, P25 comparador, P1 eliminación matemática & Huecos de relato grandes \\
4 (historia) & P18 streamgraph, P19 calendario, P14 eras, P16 títulos, P20 fiabilidad, P22 edades, P21 mapa de calor & Nueva sección «Historia» \\
5 (singulares) & P5 undercut, P2 degradación, P15 red de compañeros, P8 ghost race, P23 índice de emoción & Diferenciadores; más esfuerzo o una heurística que hay que validar \\
\end{longtable}}

\medskip\noindent\textcolor{muted!40}{\rule{\linewidth}{0.3pt}}\medskip

\section{Reorganización de la navegación}

\textbf{Situación actual}

\begin{itemize}
\item Navegación: Temporadas · Pilotos · Constructores · Récords · Calidad.
\item Pestañas de carrera: Resultado · Clasificación · Vuelta a vuelta · Neumáticos · Ritmo · Boxes · Telemetría. Son 7 y en móvil se desplazan en horizontal.
\end{itemize}

\textbf{Propuesta}

\begin{codigo}
Inicio          última carrera + mini «diferencia con el líder» (P1) + mapa|calendario (G1/P19)
Temporadas      /seasons/[year]: KPIs de relato (P11) · campeonato [puntos | diferencia | posiciones] · clasificaciones · calendario
  └ Carrera     /races/[id], 5 pestañas:
                 1. Resultado       (parrilla→V1→llegada P9 + tabla)
                 2. Clasificación   (G3 horizontal y por segmento + tabla + telemetría G10/G11 con P6 y P7 desde 2024)
                 3. Carrera         (G4 posiciones | P1b diferencias; P8 ghost race opcional)
                 4. Ritmo           (violín G8/P4 + mapa de calor P3 + tiempos G7 limitado a 2-4 pilotos)
                 5. Estrategia      (stints G5 + degradación P2 + undercuts P5 + paradas G9 + pasos por el pit lane)
Pilotos         ficha: KPIs con contexto · ADN P13 · puntos normalizables P14 · H2H divergente P12 · posiciones P21
Constructores   ficha: KPIs · temporadas (arreglo pre-1958) · alineación y motor P17 · fiabilidad del equipo P20
Circuitos       NUEVA (P24)
Comparar        NUEVA (P25 + grados de separación P15)
Historia        Récords (tabla ordenable + P14) · títulos P16 · streamgraph P18 · fiabilidad P20 · edades P22 · margen P23 · en casa P26
Datos           Calidad: enlace en el pie de página (es transparencia, no navegación principal)
\end{codigo}

\textbf{Motivos}

\begin{enumerate}
\item \textbf{La telemetría pasa a Clasificación}, porque es la vuelta de clasificación. Así desaparece el malentendido de verla en una pestaña de carrera.
\item \textbf{Neumáticos y Boxes se fusionan en «Estrategia»}: la parada y el compuesto son la misma decisión, y el detector de undercut necesita las dos cosas.
\item \textbf{«Ritmo» deja de ser un muro de 20 líneas}: el violín es la vista principal y el detalle vuelta a vuelta es secundario.
\item \textbf{5 pestañas en lugar de 7}, que caben en 375 px sin desplazamiento.
\item \textbf{Estado compartido entre pestañas}: los pilotos seleccionados en la URL (\code{?d=\allowbreak{}VER,\allowbreak{}NOR}) se mantienen en Carrera, Ritmo y Estrategia. Hoy cada gráfico tiene sus propias casillas.
\item \textbf{Récords se convierte en «Historia»}, con subsecciones. Es donde viven las vistas por eras.
\item \textbf{Circuitos y Comparar} cubren los dos huecos de entidad más evidentes. Con 7 entradas en la navegación principal hay que revisar el menú en móvil (hoy se desplaza en horizontal).
\item \textbf{Enlaces cruzados}: al pulsar una línea o barra (ECharts \code{on(\allowbreak{}'click')}) se abre la ficha del piloto, del equipo o de la carrera, y los nombres de circuito de las tablas enlazan a P24.
\end{enumerate}

\medskip\noindent\textcolor{muted!40}{\rule{\linewidth}{0.3pt}}\medskip

\section{Notas para el informe LaTeX}

\textbf{Correspondencia con el TFG}

{\small\sloppy\setlength{\tabcolsep}{5pt}\renewcommand{\arraystretch}{1.15}
\begin{longtable}{@{}>{\raggedright\arraybackslash}p{0.337\dimexpr(\linewidth-20pt)\relax}>{\raggedright\arraybackslash}p{0.663\dimexpr(\linewidth-20pt)\relax}@{}}
\toprule
\textbf{Figura del TFG} & \textbf{Elemento de la web} \\
\midrule\endhead
\bottomrule\endlastfoot
5.23 & G1 \\
5.24-5.25 & G16 y la tabla de récords \\
5.26 & KPIs de temporada y tablas \\
5.27 & G3 \\
5.28 & G4 \\
5.29 & G5 y G6 \\
5.30 & G7 \\
5.31 & G8 \\
5.32 & G9 \\
5.33 & KPIs del piloto, G12, G13 y G14 \\
\end{longtable}}

Extras sin equivalente en el TFG: G2, G10, G11, G15, la tabla de pasos por el pit lane y \code{/\allowbreak{}quality}.

\textbf{Mejoras frente al TFG que se pueden defender}

\begin{itemize}
\item Accesibilidad: resumen y tabla en cada gráfico.
\item Bilingüe.
\item Sin licencia ni cuenta de Power BI.
\item Filtro de vueltas neutralizadas más riguroso: incluye la bandera roja y la heurística del 150 \%.
\item Letras en los compuestos y texturas en los compañeros.
\item Paleta apta para daltonismo.
\item Todas las páginas funcionan sin JS salvo los gráficos, con formularios GET.
\end{itemize}

\textbf{Limitaciones heredadas del TFG que conviene reconocer}

\begin{itemize}
\item La métrica de clasificación mezcla segmentos.
\item El «tiempo de parada» es en realidad el tiempo en el pit lane.
\item Los puntos no están normalizados entre eras.
\end{itemize}

\textbf{Capturas}: tomadas el 29/09/2026 a 640 px de ancho CSS con escala de dispositivo 1,25 (PNG de 480 a 1 280 px de ancho). El recorte de la captura del navegador integrado impedía hacerlas a más ancho.
